


% ─────────────────────────────────────────────────────────────────────────────
\section{Planned Extensions}
% ─────────────────────────────────────────────────────────────────────────────

The following are flagged as future development areas; they are not yet implemented in the library, 
but they are on the roadmap. That said, no promises are made about the timeline for any of 
these features, and they may be subject to change as the project evolves.
Each will warrant its own section in this document.

\subsection{Arithmetic on Continued Fractions}

Gosper (1972) showed how to perform the four arithmetic operations on SCFs
in a streaming fashion using $2 \times 2$ and $2 \times 2 \times 2$ integer matrices,
without fully evaluating either operand.
This is the natural future direction for \texttt{FiniteSimpleContinuedFraction} arithmetic. See \cite{gosper1972}

\subsection{Generalised Continued Fractions}

The general form $b_0 + K_{n=1}^{\infty} a_n / b_n$ with arbitrary integer or
function sequences.
Applications: rapid evaluation of $\pi$, $e$, $\ln 2$, Bessel functions, etc.

\subsection{Metric Theory and Gauss Map}

The map $T: x \mapsto \{1/x\}$ (fractional part of $1/x$) on $(0,1)$ is the
\emph{Gauss map}; its invariant measure $\mu(A) = \frac{1}{\ln 2}\int_A \frac{dx}{1+x}$
governs the distribution of partial quotients (Gauss--Kuzmin statistics).
This will be relevant for statistical analysis routines.

\subsection{Multidimensional Generalisations}

Jacobi--Perron and related algorithms for simultaneous approximation; connection to
lattice reduction (LLL algorithm).